\documentclass[14pt,letterpaper]{extarticle}
\usepackage[utf8]{inputenc}

% Set margins
% https://tex.stackexchange.com/a/327136
% this should replace use of the 'fullpage' package
\usepackage[
    includehead,
    margin=2.5cm,
    bmargin=2cm,
]{geometry}

\usepackage{amssymb}
\usepackage{amsthm}
\usepackage{verbatim} % This works better than comment for some reason
\usepackage{hyperref}
\usepackage{amsmath}
\usepackage{gensymb}
\usepackage{enumitem}
\usepackage{xfrac} % Fractions - https://tex.stackexchange.com/a/230973
\usepackage{xcolor}
% The magic latex drawing thing

\usepackage{tikz}
\usetikzlibrary{shapes}
\usepackage[en-US,showdow]{datetime2}

% Colored hyperlinks:
% https://www.overleaf.com/learn/latex/Hyperlinks
\hypersetup{
    colorlinks=true,
    urlcolor=blue
}

% Setup nice headings
% https://www.overleaf.com/learn/latex/Headers_and_footers
\usepackage{fancyhdr}
\usepackage{titling} % Save title for use in fancyhdr

% Shortcuts
\newcommand*{\bP}{\mathbf{P}}



% Use black tombstone for QED
\renewcommand{\qedsymbol}{$\blacksquare$}

% Important sets
\newcommand*{\C}{\mathbb{C}}
\newcommand*{\R}{\mathbb{R}}
\newcommand*{\Q}{\mathbb{Q}}
\newcommand*{\Z}{\mathbb{Z}}

%
% symbols
%

% name by meaning, not by character
\newcommand{\grad}{\nabla}


% New command for 'partial QED' using white tombstone
\newcommand*{\partialqed}{$\hfill \square$ \\}

\newcommand*{\todo}{{\color{red}\textbf{TODO:}} }

% Increased spaces spacing
\usepackage{setspace} \onehalfspacing

\title{Math 468 - Homework 5}
\author{Nicholas Schlabach (nickninja@arizona.edu)}
\DTMsavedate{titledate}{2026-02-25}
\date{\DTMUsedate{titledate}}

%% Copied from https://tex.stackexchange.com/a/2238
%% Modified by Techcable 9/19/21 to use '[...]' instead of '(...)'
\newenvironment{amatrix}[1]{%
  \left[\begin{array}{@{}*{#1}{c}|c@{}}
}{%
  \end{array}\right]
}

\begin{document}
\pagestyle{fancy}
\fancyhead[L]{Nicholas Schlabach}
\fancyhead[C]{\hspace{5mm}\Big(\thetitle{}\Big)}
\fancyhead[R]{{\DTMsetdatestyle{mmddyy}\DTMusedate{titledate}}, Page \thepage}
\fancyfoot{} % no footer
\maketitle
% maketitle automatically sets pagestyle=plain, override this
\thispagestyle{fancy}


\begin{center}
\begin{enumerate}
\item
    \begin{comment}
    There are two ways to approach this problem.
    \vspace{3mm}\\
    Approach 1 (the way in the book):
    \end{comment}
    I am going to continue using zero-indexed entries here, as switching back and forth is too confusing. \\
    This means the communication classes are $\{0,1,2\}, \{3\}, \{4,5\}$. \\
    As demonstrated on page 49, the stationary distribution $\pi$ can be found by solving $\pi\bP=\pi$ subject to the condition $\sum \pi_i=1$.
    \begin{align*}
        \frac{\pi(0)}{2}+\frac{\pi(2)}{3} & =\pi(0) \\
        \frac{\pi(0)}{2}& =\pi(1) \\
        \pi(1) &= \pi(2)  \\
        \frac{\pi(2)}{3}+\frac{\pi(3)}{2}& =\pi(3) \\
        \frac{\pi(2)}{3} +\frac{\pi(3)}{2}+\pi(5)&=\pi(4)  \\
        \pi(4)& =\pi(5) \\
        \sum_{i=0}^5 \pi(i) & =1
    \end{align*}
    This can be written as an augmented matrix 
    \[ \begin{amatrix}{6}
        -1/2 & 0 & 1/3 & 0 & 0 & 0 & 0 \\
        1/2 & -1 & 0 & 0 & 0 & 0& 0 \\
        0 & 1 & -1 & 0 & 0 & 0 &0 \\
        0 & 0 & 1/3 & -1/2 & 0 &0 & 0 \\
        0 & 0 & 1/3 & 1/2 & -1 & 1 &0 \\
        0 & 0 & 0 & 0 & 1 & -1 & 0 \\
        1 & 1 & 1 & 1 & 1 & 1  & 1
    \end{amatrix} \]
    Converting to reduced row-echelon form gives the following:
    \[ \begin{amatrix}{6}
        1 & 0 & 0 & 0 &0 &0 & 0 \\
        0 & 1 & 0 & 0 & 0 &0 &0 \\
        0 & 0 & 1 & 0 & 0& 0& 0\\
        0 & 0 & 0 & 1 & 0 &0 &0 \\
        0 & 0 & 0 & 0 & 1 &0 &\frac{1}{2} \\
        0 & 0 & 0 & 0 & 0 &1 &\frac{1}{2} \\
        0 & 0 & 0 & 0 & 0 &0 &0 \\

    \end{amatrix} \]
    This means there is only one equilibrium solution for the system:
    \[ \pi=\langle 0,0,0,0,\frac{1}{2},\frac{1}{2}\rangle \] 
\end{enumerate}
\end{center}
\end{document}